% \documentclass{article}
% \usepackage{graphicx}
% \usepackage{amsmath} % <-- Add this line right here!

% \title{Learning Latex}
% \author{Shafnan Wasif}
% \date{June 2026}

% \begin{document}
% \maketitle
% \tableofcontents
% \section{First Demo}
% My first line in latex.

% \begin{align}
%     f(x) &= x^2 + 2x + 1 \\
%     &= (x + 1)^2
% \end{align}

% \section{Structure}
%   Utilising \emph{emphasis \emph{nested emphasis} cool}
% %emphasis surrounding text er theke differnet format e likhe
% %So, text italic mode e thakle straight korbe and vice versa
% % \textit shob kichu ke italic i korbe


% Here is an unordered list :
% \begin{itemize}
%  \item First
%  \item Second
%  \item Third
% \end{itemize}

% Here is an ordered list :
% \begin{enumerate}
%     \item First
%     \item Second
%     \item Third
% \end{enumerate}

% \subsection{Quotation}
% \begin{quote}
%     A short quoted message 
% \end{quote}

% \begin{quotation}
%     Another short quoted message but using quotation
% \end{quotation}

% % quote — an environment for a short quoted passage;
% %  it indents and adds spacing automatically. (There's a sibling,
% %   quotation, for longer multi-paragraph quotes — it adds paragraph-indentation too,
% % which quote doesn't.
% %  Lesson territory you'll bump into eventually, not now.)


% \end{document}

% %Lesson-5

% \documentclass[12pt, a4paper]{article}

% \usepackage[margin=2.5cm]{geometry}
% \usepackage{setspace}

% \onehalfspacing % calling a command from the package setspace


% \begin{document}

% \section{Layout Test}

% This paragraph should now sit inside 2.5cm margins, on A4 paper,
% at 12pt base font size, with one-and-a-half line spacing instead
% of the default single spacing.

% This is a second paragraph, so you can see the line spacing
% clearly across multiple lines of text rather than just one.

% \end{document}

% % Lesson-6

% \documentclass{article}

% \usepackage{amsmath}
% \usepackage{amssymb} % <-- ADD THIS LINE
% \usepackage{xcolor}

% \newcommand{\R}{\mathbb{R}}
% \newcommand{\highlight}[1]{\textcolor{red}{#1}}

% \begin{document}

% \section{Packages and Custom Commands}

% The amsmath package gives access to better math environments,
% which we'll use properly in the math lesson. For now, just note
% that $x \in \R$ uses our custom shortcut instead of typing
% \verb|\mathbb{R}| every time.

% Here is some \highlight{highlighted text} using a command we
% defined ourselves, rather than one a package gave us.

% \end{document}

% % The syntax is \newcommand{\name}{replacement text} — anywhere you later type \name, 
% % LaTeX substitutes the replacement text verbatim, as if you'd typed it there yourself.

% % \newcommand{\customcolor}[2]{\textcolor{#1}{#2}} -- takes in 2 arguments 
% %  This is \customcolor{blue}{sky blue} and this is \customcolor{green}{forest green}.



% Figure environment is a float . they float to the best possible position.

\documentclass{article}
\usepackage{graphicx}
\usepackage{multirow}
\usepackage{booktabs}
\begin{document}

\section{Working With Images}

Here is a paragraph of text before the figure. The figure below
is not guaranteed to appear exactly here -- LaTeX may move it to
wherever it fits best on the page.

\begin{figure}[h] % [h] =="Please try to place this figure exactly where it appears in the source text, if there is enough room on the page."
%  b-bottom , t-top, p-places the float on a separate page reserved for float . 
 
  \centering
  \includegraphics[width=0.6\textwidth]{example-image.png}
  \caption{An example figure with a caption.}
  \label{fig:example}
\end{figure}

Here is a paragraph of text after the figure.



\begin{figure}
  \centering
  \includegraphics[width=0.6\textwidth]{example-image.png}
  \caption{Second example}
  \label{fig: example 2}
\end{figure}



\begin{table}[h]
  \centering
  \begin{tabular}{|l|c|c|}
    \hline
    Method & Accuracy & Time Complexity \\
    \hline
    Rasterization & Medium & $O(n)$ \\
    \hline
    Ray Tracing & High & $O(n^2)$ \\
    
     & Very High & $O(n^3)$ \\
     \hline
    \multicolumn{2}{|c|}{Hybrid Method}& $O(n^2)$ \\
    \hline
  \end{tabular}
  \caption{Comparison of Image Synthesis Techniques}
  \label{tab:comparison}
\end{table}



\begin{table}[h]
  \centering
  \begin{tabular}{|l|c|c|}
    \hline
    \multirow{2}{*}{Module} & \multicolumn{2}{c|}{Score} \\
    \cline{2-3}
     & Theory & Lab \\
    \hline
    Module A & 75 & 80 \\
    Module B & 88 & 90 \\
     & 85 & 87 \\
    Module C & 70 & 72 \\
    \hline
  \end{tabular}
  \caption{Module-wise Performance Summary}
  \label{tab:modules}
\end{table}


\begin{table}[h]
\centering
  \begin{tabular}{|l|c|c|}
    \hline
    Method & Accuracy & Time Complexity \\
    \hline
    Rasterization & Medium & $O(n)$ \\
    \hline
    \multirow{2}{*}{Ray Tracing} & High & $O(n^2)$ \\
    & Very High& $O(n^3)$\\
    \hline
    \multicolumn{2}{|c|}{Hybrid Method} & $O(n^2)$ \\
    \hline
  \end{tabular}
  \caption{ Comparison of Image syntehesis techniques}
\end{table}

% Self-Practice
\begin{table}
\centering
\begin{tabular}{|c|c|c|}
  \hline
  \multirow{2}{*}{Method} & \multicolumn{2}{c|}{Score} \\
  \cline{2-3}
  & Theory & Lab \\
  \hline
  Module A & 75 & 80\\
  \multirow{2}{*}{Module B} & 88 & 90\\
  & 85 & 87\\
  \hline
  Module C & 70 & 72\\
  \hline
  
\end{tabular}
\caption{Module-wise Performance Summary}
\end{table}


% Table-5
\begin{table}[h]
  \centering
  \begin{tabular}{lcr}
    \toprule
    Item & Quantity & Price \\
    \midrule
    Apples & 10 & \$5.00 \\
    Bananas & 6 & \$2.40 \\
    Cherries & 200 & \$8.75 \\
    \bottomrule
  \end{tabular}
  \caption{Same data, booktabs style}
\end{table}


\end{document}




